Here, we provide a self-contained tutorial for the model and fitting procedure, with the goal of making it easier for readers to understand the model and reproduce the results.
The fitting procedure is from \citet{hahn2024unifying}, and we refer to the formal description there.
Here, we intend to give a user-friendly exposition to readers interested in applying the fitting procedure to psychophysical data.

The codebase provides instructions for running on a simulated dataset\footnote{\url{https://github.com/m-hahn/identifiability-bayesian-models/tree/main/code}}, on the original dataset by \citet{deGardelle2010AnOI}\footnote{\url{https://github.com/m-hahn/identifiability-bayesian-models/tree/main/code/Gardelle}} and on the dataset by \citet{Remington2018LateBI}\footnote{\url{https://gitlab.com/m-hahn/unifying-theory-biases/-/tree/main/code/Remington} in the codebase from \citet{hahn2024unifying}.}, including for visualizing the resulting fits.

We first discuss which aspects are specified by the user, and which parts are recovered by the fitting procedure.
The user specifies the following parts:
\begin{enumerate}
\item The stimulus space $\mathcal{X}$; it may be a finite interval or a circle.
The  sensory space $\mathcal{Y}$ has the same topology; in practice, it can be chosen to be $[0,1]$ (interval) or $[0,2\pi)$ (circular).
\item The exponent $p \geq 0$ of the $L_p$ loss function. The fitting procedure admits any nonnegative integer; we use $p=0,1,2,4,6,8,\dots$.
\item The number $L$ of sensory noise levels.
\end{enumerate}
The behavioral data is expected to consist of trials that each consist of
\begin{enumerate}
\item the true stimulus $\theta \in \mathcal{X}$,
\item the observed response $x \in \mathcal{X}$,
\item an index for the sensory noise level, $l \in \{1, \dots, L\}$ (e.g., an index for a condition under a manipulation of exposure duration, contrast, or another manipulation affecting the magnitude of sensory noise).
\end{enumerate}
Given these, the fitting procedure determines
\begin{enumerate}
	\item A prior $\Prior(\theta)$, which is a probability density on $\mathcal{X}$.
	\item An encoding function $F(\theta) : \mathcal{X} \rightarrow \mathcal{Y}$; equivalently, a function $F'(\theta) : \mathcal{X} \mapsto (0,\infty)$ describing the slope of $F(\theta)$.
	\item Sensory noise magnitudes  $\sigma_1, \dots, \sigma_L > 0$, which specify the magnitude of sensory noise in the $L$ different noise level conditions
        \item The variance  $\rho^2$ of motor noise. 
	\item $\xi \in [0,1)$, the guessing rate.
\end{enumerate}
There are several possible extensions, such as to stimulus (external) noise (Section~\ref{sec:stimulus-noise}), to subject-specific adjustments (SI Appendix S2.1.4 in \citet{hahn2024unifying}), and to encoding functions varying with noise conditions (Section~\ref{sec:vary-with-noise}).

For any dataset, we use 90\% of trials for model fitting (training set) and 10\% of trials for evaluating the quality of fit (held-out or test set).
The quality-of-fit statistic is negative log-likelihood (NLL) on held-out data, which is the sum of the negative log-likelihoods of all held-out trials.

Note that the model is fitted separately for each potential loss function exponent $p$ (giving a different heldout NLL for each $p$), whereas the other parameters (prior, encoding, noise and guessing parameters) are jointly fitted.


By Bayes' rule, the posterior over $\theta$ given an encoding $m$ in the $l$-th noise level ($l = 1, \dots, L$) is
\begin{equation}
P(\theta|m; \sigma_l) \propto P_{sensory}(m|\theta; \sigma_l^2) \Prior(\theta)
\end{equation}
where $P_{sensory}(m|\theta; \sigma_l^2)$ is the sensory noise distribution over $\mathcal{Y}$ (Gaussian with variance $\sigma_l^2$ or von Mises with concentration $1/\sigma_l^2$, for the corresponding noise level $l$).
The $L_p$-estimator $\hat{\theta}$ is obtained from the posterior $P(\theta|m; \sigma_l^2)$ by minimizing the $L_p$ loss function.
Then, the data log-likelihood for a single pair of true stimulus $\theta \in \mathcal{X}$ and observed response $x \in \mathcal{X}$ at the $l$-th noise level ($l=1, \dots, L$) is defined as follows:
\begin{equation}\label{eq:data-likelihood-raw}
\log P(x|\theta; \sigma_l^2) = \log \left\{ \frac{\xi}{\operatorname{Vol}(\mathcal{X})} + (1-\xi)  \int_{m \in \mathcal{Y}} P_{motor}(x|\hat{\theta}_m) P_{sensory}(m|\theta; \sigma_l^2) dm \right\}
\end{equation}
where $\hat{\theta}_m \in \mathcal{X}$ is the $L_p$-estimator given the posterior $P(\hat{\theta}|m; \sigma_l^2)$ for an encoding $m \in \mathcal{Y}$, and $P_{motor}(x|\hat{\theta}_m)$ is the motor noise distribution over $\mathcal{X}$ (Gaussian with variance $\rho^2$ or von Mises with concentration $1/\rho^2$).
This can equivalently be expressed, by transforming the integral over $m$ to an integral over $F^{-1}(m)$, as
\begin{equation}\label{eq:data-likelihood}
\log P(x|\theta; \sigma_l^2) = \log \left\{ \frac{\xi}{\operatorname{Vol}(\mathcal{X})} + (1-\xi)  \int_{m \in \mathcal{X}} P_{motor}(x|\hat{\theta}_m) P_{sensory-transformed}(m|\theta; \sigma_l^2) dm \right\}
\end{equation}
where, by the transformation law for probability densities,
\begin{equation}
P_{sensory-transformed}(m|\theta; \sigma_l^2) \propto P_{sensory}(F(m)|\theta; \sigma_l^2) \cdot F'(m)
\end{equation}
This is advantageous for implementation because it avoids the need to separately represent $\mathcal{X}$ and $\mathcal{Y}$.

\paragraph*{Parameterization of Prior and Encoding}
The model is implemented computationally on a regularly spaced discretized grid $\theta_1, \dots, \theta_N$.
We parameterize the encoding and prior by assigning one value of $F'$, $\Prior$ to each grid point $\theta_i$.
The encoding and prior are parameterized in terms of logarithms of these values, and obtained by normalizing the exponentiated parameters (softmax transform), as follows:
\begin{align}\label{eq:fprime-param}
	F'(\theta_i) =& \operatorname{Vol}(\mathcal{Y}) \frac{\exp(\alpha(\theta_i))}{\sum_{j=1}^N\exp(\alpha(\theta_j))} \\
\label{eq:prior-param}
		\Prior(\theta_i) =& \frac{\exp(\beta(\theta_i))}{\sum_{j=1}^N\exp(\beta(\theta_j))}
\end{align}
for $\theta_i$ on the grid; here, $\alpha(\theta_i)$ and $\beta(\theta_i)$ ($i=1,\dots, N$) are real-valued parameters of the model. 
We note that this is a discretized version of the theoretical construction of a measure over models (Section~\ref{sec:measure-mu}).

\paragraph*{Implementation of Model}
We map the  input stimulus $\theta$  to the closest grid point $\theta_i$ ($i = 1,\dots, N$) and then compute likelihoods and posteriors exactly 
on the discretized grid, and obtain the $L_p$ estimator via Newton's method.
First, the sensory likelihood is, for the $l$-th noise level ($k, i = 1, \dots, N$):
\begin{equation}
P_{sensory-transformed}(\theta_k|\theta_i; \sigma_l^2) \propto \exp(-\frac{(F(\theta_k) - F(\theta_i))^2}{2\sigma_l^2}) \cdot F'(\theta_k)
\end{equation}
(analogously in the case of von Mises noise).
Second, the posterior is ($j, k = 1, \dots, N$):
\begin{equation}
P(\theta_j|\theta_k; \sigma_l^2) \propto P_{sensory-transformed}(\theta_k|\theta_j; \sigma_l^2) \Prior(\theta_j)
\end{equation}
The estimator $\hat{\theta} \in \mathcal{X}$ is obtained from the posterior $P(\theta_k|\theta_i; \sigma_l^2)$ by minimizing the expected $L_p$ loss function:
\begin{equation}\label{eq:lp-estimator-discretized}
\hat{\theta} = \argmin_{\hat{\theta} \in \mathcal{X}} \sum_{k=1}^N P(\theta_k|\theta_i; \sigma_l^2) |\hat{\theta}-\theta_k|^p
\end{equation}
or, in the case of a circular space,
\begin{equation}\label{eq:lp-estimator-discretized-circular}
\hat{\theta} = \argmin_{\hat{\theta} \in \mathcal{X}} \sum_{k=1}^N P(\theta_k|\theta_i; \sigma_l^2) (1-\cos(\hat\theta-\theta))^{p/2}
\end{equation}
These formulas are used for $p \geq 2$; implementations for $p=0,1$ are described in Section~\ref{sec:implementation}. 
%
%On interval spaces, the $L_p$ loss function is defined as $|\hat\theta-\theta|^p$ (for $p>0$).
%On circular spaces, the $L_p$ loss function is defined as $(1-\cos(\hat\theta-\theta))^{p/2}$ (for $p>1$).
%The $L_0$ estimator is operationalized as the maximum of a kernel-smoothed version of the posterior (see \citet{hahn2024unifying} for details).

\paragraph*{Model Fitting}
The data log-likelihood for a single pair consisting of a true stimulus mapped to a grid point $\theta_i$ ($i=1, \dots, N$) and an observed response $x \in \mathcal{X}$, under the $l$-th noise level ($l=1,\dots,L$), is defined by discretizing the integral in Eq.~\eqref{eq:data-likelihood}:
\begin{equation}\label{eq:data-likelihood-discretized}
\log P(x|\theta_i; \sigma_l^2) = \log \left\{\frac{\xi}{\operatorname{Vol}(\mathcal{X})} + (1-\xi)  \sum_k P_{motor}(x|\hat{\theta}_{\theta_k}) P_{sensory-transformed}(\theta_k|\theta_i; \sigma_l^2) \right\}
\end{equation}
where $k$ runs over the points on the discretized grid, and $\hat{\theta}_{\theta_k}$ is the $L_p$ estimator derived from posterior $P(\theta|\theta_k; \sigma_l^2)$.
Summing this over all datapoints gives the log-likelihood of the behavioral dataset.

We fit all model parameters (prior, encoding, noise magnitudes, guessing rate) end-to-end via gradient descent to maximize the likelihood of the behavioral data.
We implement (\ref{eq:data-likelihood-discretized}) in PyTorch \citep{DBLP:conf/nips/PaszkeGMLBCKLGA19}, which provides differentiation out-of-the box.
Notably, we provide a custom implementation for the $L_p$ optimizers as standalone PyTorch modules, allowing differentiating through $\hat{\theta}$. In the forward pass, this module obtains $\hat{\theta}$ from the posterior $P(\theta|\theta_k; \sigma_l^2)$ by minimizing (\ref{eq:lp-estimator-discretized}--\ref{eq:lp-estimator-discretized-circular}) via Newton's method. For the backward pass, we compute the derivatives of $\hat{\theta}$ w.r.t. the weights of the posterior $P(\theta_j|\theta_k; \sigma_l^2)$ ($j, k = 1, \dots, N$) via implicit differentiation.
We refer to \citet{hahn2024unifying} for technical details.

\citet{hahn2024unifying} mitigate overfitting in the shapes of prior and encoding by adding a regularization term penalizing\footnote{For circular spaces, $\left(\alpha(\theta_N), \alpha(\theta_{1})\right)^2$ is added additionally.} $\sum_{i=1}^{N-1} \left(\alpha(\theta_i), \alpha(\theta_{i+1})\right)^2$ and similarly for $\beta$, modulated with a regularization strength $\lambda$ that is set by the user.
In principle, one can determine the regularization strength for a given behavioral dataset via cross-validation, though this comes with computational cost. 
As described in Materials and Methods, we used the same regularization strength as was identified for the data from \citet{deGardelle2010AnOI} by \citet{hahn2024unifying}, and used this across all datasets.
A control study (Section~\ref{sec:robustness-reg}) shows that the regularization strength has little impact on the identifiability.
